
\section{Long-term and Short-term Forecasting}\label{appx:long-short-term}

\subsection{Long-term Forecasting}

By leveraging two-stage \method{} framework while bypassing pre-alignment, our method attains SOTA performance in \textbf{35} instances across eight time series benchmarks. This underscores the considerable potential of LLMs as robust and reliable time series forecasters.

\shortautoref{tab:long-term-forecasting-full} presents a comprehensive evaluation of long-term forecasting performance across eight benchmark datasets, categorizing baseline models by their input modality (Language, Multi-Modal, Visual, and Numerical).

\begin{table}[H]\centering\small
\caption{Full long-term forecasting results. Each entry is MSE / MAE; lower values are better.}\label{tab:long-term-forecasting-full}
\textbf{ETTh1}\par\smallskip
\begin{adjustbox}{max width=\linewidth}\begin{tabular}{lrrrrr}\toprule
\textbf{Model} & \textbf{96} & \textbf{192} & \textbf{336} & \textbf{720} & \textbf{Avg}\\\midrule
\rowcolor{blue!5}
\textbf{RDTU} & \textbf{0.351 / 0.382} & \textbf{0.389 / 0.403} & \textbf{0.412 / 0.419} & \textbf{0.442 / 0.446} & \textbf{0.399 / 0.413}\\
GPT4TS & 0.370 / 0.389 & 0.412 / 0.413 & 0.448 / 0.431 & 0.441 / 0.449 & 0.418 / 0.421\\
Time-LLM & 0.376 / 0.402 & 0.407 / 0.421 & 0.430 / 0.438 & 0.457 / 0.468 & 0.418 / 0.432\\
DMMV-A & 0.354 / 0.389 & 0.393 / 0.405 & 0.387 / 0.413 & 0.445 / 0.450 & 0.395 / 0.414\\
Time-VLM & 0.361 / 0.386 & 0.397 / 0.415 & 0.420 / 0.421 & 0.441 / 0.458 & 0.405 / 0.420\\
VisionTS & 0.355 / 0.386 & 0.395 / 0.407 & 0.419 / 0.421 & 0.458 / 0.460 & 0.407 / 0.419\\
PatchTST & 0.370 / 0.399 & 0.413 / 0.421 & 0.422 / 0.436 & 0.447 / 0.466 & 0.413 / 0.431\\
CycleNet & 0.374 / 0.396 & 0.406 / 0.415 & 0.431 / 0.430 & 0.450 / 0.464 & 0.415 / 0.426\\
TimesNet & 0.384 / 0.402 & 0.436 / 0.429 & 0.491 / 0.469 & 0.521 / 0.500 & 0.458 / 0.450\\
DLinear & 0.375 / 0.399 & 0.405 / 0.416 & 0.439 / 0.416 & 0.472 / 0.490 & 0.423 / 0.430\\
FEDformer & 0.376 / 0.419 & 0.420 / 0.448 & 0.459 / 0.465 & 0.506 / 0.507 & 0.440 / 0.460\\
\bottomrule\end{tabular}\end{adjustbox}\end{table}
\begin{table}[H]\centering\small
\caption*{Table S10 (continued). ETTh2}
\textbf{ETTh2}\par\smallskip
\begin{adjustbox}{max width=\linewidth}\begin{tabular}{lrrrrr}\toprule
\textbf{Model} & \textbf{96} & \textbf{192} & \textbf{336} & \textbf{720} & \textbf{Avg}\\\midrule
\rowcolor{blue!5}
\textbf{RDTU} & \textbf{0.263 / 0.329} & \textbf{0.334 / 0.372} & \textbf{0.355 / 0.381} & \textbf{0.395 / 0.425} & \textbf{0.337 / 0.377}\\
GPT4TS & 0.280 / 0.335 & 0.348 / 0.380 & 0.380 / 0.405 & 0.406 / 0.436 & 0.354 / 0.389\\
Time-LLM & 0.286 / 0.346 & 0.361 / 0.391 & 0.390 / 0.414 & 0.405 / 0.434 & 0.361 / 0.396\\
DMMV-A & 0.294 / 0.349 & 0.339 / 0.395 & 0.322 / 0.384 & 0.392 / 0.425 & 0.337 / 0.388\\
Time-VLM & 0.267 / 0.335 & 0.326 / 0.373 & 0.357 / 0.406 & 0.412 / 0.449 & 0.341 / 0.391\\
VisionTS & 0.288 / 0.334 & 0.349 / 0.380 & 0.364 / 0.398 & 0.403 / 0.431 & 0.351 / 0.386\\
PatchTST & 0.274 / 0.336 & 0.339 / 0.379 & 0.329 / 0.380 & 0.379 / 0.422 & 0.330 / 0.379\\
CycleNet & 0.279 / 0.341 & 0.342 / 0.385 & 0.371 / 0.413 & 0.426 / 0.451 & 0.355 / 0.398\\
TimesNet & 0.340 / 0.374 & 0.402 / 0.414 & 0.452 / 0.452 & 0.462 / 0.468 & 0.414 / 0.427\\
DLinear & 0.289 / 0.353 & 0.383 / 0.418 & 0.448 / 0.465 & 0.605 / 0.551 & 0.431 / 0.447\\
FEDformer & 0.358 / 0.397 & 0.429 / 0.439 & 0.496 / 0.487 & 0.463 / 0.474 & 0.437 / 0.449\\
\bottomrule\end{tabular}\end{adjustbox}\end{table}
\begin{table}[H]\centering\small
\caption*{Table S10 (continued). ETTm1}
\textbf{ETTm1}\par\smallskip
\begin{adjustbox}{max width=\linewidth}\begin{tabular}{lrrrrr}\toprule
\textbf{Model} & \textbf{96} & \textbf{192} & \textbf{336} & \textbf{720} & \textbf{Avg}\\\midrule
\rowcolor{blue!5}
\textbf{RDTU} & \textbf{0.272 / 0.326} & \textbf{0.327 / 0.354} & \textbf{0.350 / 0.375} & \textbf{0.410 / 0.412} & \textbf{0.340 / 0.367}\\
GPT4TS & 0.300 / 0.340 & 0.343 / 0.368 & 0.376 / 0.386 & 0.431 / 0.416 & 0.363 / 0.378\\
Time-LLM & 0.291 / 0.341 & 0.341 / 0.369 & 0.359 / 0.379 & 0.433 / 0.419 & 0.356 / 0.377\\
DMMV-A & 0.279 / 0.329 & 0.317 / 0.357 & 0.351 / 0.381 & 0.411 / 0.415 & 0.340 / 0.371\\
Time-VLM & 0.304 / 0.346 & 0.332 / 0.366 & 0.364 / 0.383 & 0.402 / 0.410 & 0.351 / 0.376\\
VisionTS & 0.284 / 0.332 & 0.327 / 0.362 & 0.354 / 0.382 & 0.411 / 0.415 & 0.344 / 0.373\\
PatchTST & 0.290 / 0.342 & 0.332 / 0.369 & 0.366 / 0.392 & 0.416 / 0.420 & 0.351 / 0.381\\
CycleNet & 0.299 / 0.348 & 0.334 / 0.367 & 0.368 / 0.386 & 0.417 / 0.414 & 0.355 / 0.379\\
TimesNet & 0.338 / 0.375 & 0.374 / 0.387 & 0.410 / 0.411 & 0.478 / 0.450 & 0.400 / 0.406\\
DLinear & 0.299 / 0.343 & 0.335 / 0.365 & 0.369 / 0.386 & 0.425 / 0.421 & 0.357 / 0.379\\
FEDformer & 0.379 / 0.419 & 0.426 / 0.441 & 0.445 / 0.459 & 0.543 / 0.490 & 0.448 / 0.452\\
\bottomrule\end{tabular}\end{adjustbox}\end{table}
\begin{table}[H]\centering\small
\caption*{Table S10 (continued). ETTm2}
\textbf{ETTm2}\par\smallskip
\begin{adjustbox}{max width=\linewidth}\begin{tabular}{lrrrrr}\toprule
\textbf{Model} & \textbf{96} & \textbf{192} & \textbf{336} & \textbf{720} & \textbf{Avg}\\\midrule
\rowcolor{blue!5}
\textbf{RDTU} & \textbf{0.157 / 0.244} & \textbf{0.219 / 0.288} & \textbf{0.270 / 0.323} & \textbf{0.354 / 0.379} & \textbf{0.250 / 0.309}\\
GPT4TS & 0.163 / 0.249 & 0.222 / 0.291 & 0.273 / 0.327 & 0.357 / 0.376 & 0.254 / 0.311\\
Time-LLM & 0.162 / 0.248 & 0.235 / 0.304 & 0.280 / 0.329 & 0.366 / 0.382 & 0.261 / 0.316\\
DMMV-A & 0.172 / 0.260 & 0.227 / 0.298 & 0.272 / 0.327 & 0.351 / 0.381 & 0.256 / 0.317\\
Time-VLM & 0.160 / 0.250 & 0.215 / 0.291 & 0.270 / 0.325 & 0.348 / 0.378 & 0.248 / 0.311\\
VisionTS & 0.174 / 0.262 & 0.228 / 0.297 & 0.281 / 0.337 & 0.384 / 0.410 & 0.267 / 0.327\\
PatchTST & 0.165 / 0.255 & 0.220 / 0.292 & 0.274 / 0.329 & 0.362 / 0.385 & 0.255 / 0.315\\
CycleNet & 0.159 / 0.247 & 0.214 / 0.286 & 0.269 / 0.322 & 0.363 / 0.382 & 0.251 / 0.309\\
TimesNet & 0.187 / 0.267 & 0.249 / 0.309 & 0.321 / 0.351 & 0.408 / 0.403 & 0.291 / 0.333\\
DLinear & 0.167 / 0.260 & 0.224 / 0.303 & 0.281 / 0.342 & 0.397 / 0.421 & 0.267 / 0.332\\
FEDformer & 0.203 / 0.287 & 0.269 / 0.328 & 0.325 / 0.366 & 0.421 / 0.415 & 0.305 / 0.349\\
\bottomrule\end{tabular}\end{adjustbox}\end{table}
\begin{table}[H]\centering\small
\caption*{Table S10 (continued). Illness}
\textbf{Illness}\par\smallskip
\begin{adjustbox}{max width=\linewidth}\begin{tabular}{lrrrrr}\toprule
\textbf{Model} & \textbf{24} & \textbf{36} & \textbf{48} & \textbf{60} & \textbf{Avg}\\\midrule
\rowcolor{blue!5}
\textbf{RDTU} & \textbf{1.609 / 0.771} & \textbf{1.453 / 0.835} & \textbf{1.545 / 0.841} & \textbf{1.630 / 0.802} & \textbf{1.559 / 0.812}\\
GPT4TS & 1.869 / 0.823 & 1.853 / 0.854 & 1.886 / 0.855 & 1.877 / 0.877 & 1.871 / 0.852\\
Time-LLM & 1.792 / 0.807 & 1.833 / 0.833 & 2.269 / 1.012 & 2.177 / 0.925 & 2.018 / 0.894\\
DMMV-A & 1.409 / 0.754 & 1.290 / 0.745 & 1.499 / 0.810 & 1.428 / 0.773 & 1.407 / 0.771\\
Time-VLM & -- & -- & -- & -- & --\\
VisionTS & 1.613 / 0.834 & 1.316 / 0.750 & 1.548 / 0.818 & 1.450 / 0.783 & 1.482 / 0.796\\
PatchTST & 1.319 / 0.754 & 1.430 / 0.834 & 1.553 / 0.815 & 1.470 / 0.788 & 1.443 / 0.798\\
CycleNet & 2.255 / 1.017 & 2.121 / 0.950 & 2.187 / 1.007 & 2.185 / 0.997 & 2.187 / 0.992\\
TimesNet & 2.317 / 0.934 & 1.972 / 0.920 & 2.238 / 0.940 & 2.027 / 0.928 & 2.139 / 0.931\\
DLinear & 2.215 / 1.081 & 1.963 / 0.963 & 2.130 / 1.024 & 2.368 / 1.096 & 2.169 / 1.041\\
FEDformer & 3.228 / 1.260 & 2.679 / 1.080 & 2.622 / 1.078 & 2.857 / 1.157 & 2.847 / 1.144\\
\bottomrule\end{tabular}\end{adjustbox}\end{table}
\begin{table}[H]\centering\small
\caption*{Table S10 (continued). Electricity}
\textbf{Electricity}\par\smallskip
\begin{adjustbox}{max width=\linewidth}\begin{tabular}{lrrrrr}\toprule
\textbf{Model} & \textbf{96} & \textbf{192} & \textbf{336} & \textbf{720} & \textbf{Avg}\\\midrule
\rowcolor{blue!5}
\textbf{RDTU} & \textbf{0.124 / 0.210} & \textbf{0.144 / 0.235} & \textbf{0.162 / 0.257} & \textbf{0.193 / 0.281} & \textbf{0.156 / 0.246}\\
GPT4TS & 0.141 / 0.239 & 0.158 / 0.253 & 0.172 / 0.266 & 0.207 / 0.293 & 0.170 / 0.263\\
Time-LLM & 0.137 / 0.233 & 0.152 / 0.247 & 0.169 / 0.267 & 0.200 / 0.290 & 0.165 / 0.259\\
DMMV-A & 0.126 / 0.213 & 0.145 / 0.237 & 0.162 / 0.254 & 0.197 / 0.286 & 0.158 / 0.248\\
Time-VLM & 0.142 / 0.245 & 0.157 / 0.260 & 0.174 / 0.276 & 0.214 / 0.308 & 0.172 / 0.272\\
VisionTS & 0.127 / 0.217 & 0.148 / 0.237 & 0.163 / 0.253 & 0.199 / 0.293 & 0.159 / 0.250\\
PatchTST & 0.129 / 0.222 & 0.157 / 0.240 & 0.163 / 0.259 & 0.197 / 0.290 & 0.162 / 0.253\\
CycleNet & 0.128 / 0.223 & 0.144 / 0.237 & 0.160 / 0.254 & 0.198 / 0.287 & 0.158 / 0.250\\
TimesNet & 0.168 / 0.272 & 0.184 / 0.289 & 0.198 / 0.300 & 0.220 / 0.320 & 0.193 / 0.295\\
DLinear & 0.140 / 0.237 & 0.153 / 0.249 & 0.169 / 0.267 & 0.203 / 0.301 & 0.166 / 0.264\\
FEDformer & 0.193 / 0.308 & 0.201 / 0.315 & 0.214 / 0.329 & 0.246 / 0.355 & 0.214 / 0.327\\
\bottomrule\end{tabular}\end{adjustbox}\end{table}
\begin{table}[H]\centering\small
\caption*{Table S10 (continued). Weather}
\textbf{Weather}\par\smallskip
\begin{adjustbox}{max width=\linewidth}\begin{tabular}{lrrrrr}\toprule
\textbf{Model} & \textbf{96} & \textbf{192} & \textbf{336} & \textbf{720} & \textbf{Avg}\\\midrule
\rowcolor{blue!5}
\textbf{RDTU} & \textbf{0.141 / 0.184} & \textbf{0.189 / 0.227} & \textbf{0.239 / 0.272} & \textbf{0.314 / 0.322} & \textbf{0.221 / 0.251}\\
GPT4TS & 0.148 / 0.188 & 0.192 / 0.230 & 0.246 / 0.273 & 0.320 / 0.328 & 0.227 / 0.255\\
Time-LLM & 0.155 / 0.199 & 0.223 / 0.261 & 0.251 / 0.279 & 0.345 / 0.342 & 0.244 / 0.270\\
DMMV-A & 0.143 / 0.195 & 0.187 / 0.242 & 0.237 / 0.273 & 0.302 / 0.315 & 0.217 / 0.256\\
Time-VLM & 0.148 / 0.200 & 0.193 / 0.240 & 0.243 / 0.281 & 0.312 / 0.332 & 0.224 / 0.263\\
VisionTS & 0.146 / 0.191 & 0.194 / 0.238 & 0.243 / 0.275 & 0.318 / 0.328 & 0.225 / 0.258\\
PatchTST & 0.149 / 0.198 & 0.194 / 0.241 & 0.245 / 0.282 & 0.314 / 0.334 & 0.226 / 0.264\\
CycleNet & 0.167 / 0.221 & 0.212 / 0.258 & 0.260 / 0.293 & 0.328 / 0.339 & 0.242 / 0.278\\
TimesNet & 0.172 / 0.220 & 0.219 / 0.261 & 0.280 / 0.306 & 0.365 / 0.359 & 0.259 / 0.287\\
DLinear & 0.176 / 0.237 & 0.220 / 0.282 & 0.265 / 0.319 & 0.333 / 0.362 & 0.249 / 0.300\\
FEDformer & 0.217 / 0.296 & 0.276 / 0.336 & 0.339 / 0.380 & 0.403 / 0.428 & 0.309 / 0.360\\
\bottomrule\end{tabular}\end{adjustbox}\end{table}
\begin{table}[H]\centering\small
\caption*{Table S10 (continued). Traffic}
\textbf{Traffic}\par\smallskip
\begin{adjustbox}{max width=\linewidth}\begin{tabular}{lrrrrr}\toprule
\textbf{Model} & \textbf{96} & \textbf{192} & \textbf{336} & \textbf{720} & \textbf{Avg}\\\midrule
\rowcolor{blue!5}
\textbf{RDTU} & \textbf{0.358 / 0.235} & \textbf{0.375 / 0.248} & \textbf{0.388 / 0.256} & \textbf{0.435 / 0.280} & \textbf{0.389 / 0.255}\\
GPT4TS & 0.396 / 0.264 & 0.412 / 0.268 & 0.421 / 0.273 & 0.455 / 0.291 & 0.421 / 0.274\\
Time-LLM & 0.392 / 0.267 & 0.409 / 0.271 & 0.434 / 0.296 & 0.451 / 0.291 & 0.422 / 0.281\\
DMMV-A & 0.344 / 0.237 & 0.363 / 0.249 & 0.387 / 0.256 & 0.433 / 0.284 & 0.389 / 0.257\\
Time-VLM & 0.393 / 0.290 & 0.405 / 0.296 & 0.420 / 0.305 & 0.459 / 0.323 & 0.419 / 0.304\\
VisionTS & 0.346 / 0.232 & 0.376 / 0.245 & 0.389 / 0.252 & 0.432 / 0.293 & 0.386 / 0.256\\
PatchTST & 0.360 / 0.249 & 0.379 / 0.256 & 0.392 / 0.264 & 0.432 / 0.286 & 0.391 / 0.264\\
CycleNet & 0.397 / 0.278 & 0.411 / 0.283 & 0.424 / 0.289 & 0.450 / 0.305 & 0.421 / 0.289\\
TimesNet & 0.593 / 0.321 & 0.617 / 0.336 & 0.629 / 0.336 & 0.640 / 0.350 & 0.620 / 0.336\\
DLinear & 0.410 / 0.282 & 0.423 / 0.287 & 0.436 / 0.296 & 0.466 / 0.315 & 0.434 / 0.295\\
FEDformer & 0.587 / 0.366 & 0.604 / 0.373 & 0.621 / 0.383 & 0.626 / 0.382 & 0.610 / 0.376\\
\bottomrule\end{tabular}\end{adjustbox}\end{table}
\noindent\textbf{Source-reported ranking counts.} RDTU: 35, GPT4TS: 2, Time-LLM: 0, DMMV-A: 23, Time-VLM: 6, VisionTS: 6, PatchTST: 7, CycleNet: 6, TimesNet: 0, DLinear: 0, FEDformer: 0. The original counting convention is not specified; these values are retained without recalculation.\par


\begin{table}[H]\centering\small
\caption{Long-term forecasting: additional baselines. Each entry is MSE / MAE; lower values are better.}\label{tab:long-term-forecasting-additional}
\textbf{ETTh1}\par\smallskip
\begin{adjustbox}{max width=\linewidth}\begin{tabular}{lrrrrr}\toprule
\textbf{Model} & \textbf{96} & \textbf{192} & \textbf{336} & \textbf{720} & \textbf{Avg}\\\midrule
\rowcolor{blue!5}
\textbf{RDTU} & \textbf{0.351 / 0.382} & \textbf{0.389 / 0.403} & \textbf{0.412 / 0.419} & \textbf{0.442 / 0.446} & \textbf{0.399 / 0.413}\\
Autoformer & 0.449 / 0.459 & 0.500 / 0.482 & 0.521 / 0.496 & 0.514 / 0.512 & 0.496 / 0.487\\
Stationary & 0.513 / 0.491 & 0.534 / 0.504 & 0.588 / 0.535 & 0.643 / 0.616 & 0.570 / 0.537\\
ETSformer & 0.494 / 0.479 & 0.538 / 0.504 & 0.574 / 0.521 & 0.562 / 0.535 & 0.542 / 0.510\\
LightTS & 0.424 / 0.432 & 0.475 / 0.462 & 0.518 / 0.488 & 0.547 / 0.533 & 0.491 / 0.479\\
Informer & 0.865 / 0.713 & 1.008 / 0.792 & 1.107 / 0.809 & 1.181 / 0.865 & 1.040 / 0.795\\
Reformer & 0.837 / 0.728 & 0.923 / 0.766 & 1.097 / 0.835 & 1.257 / 0.889 & 1.029 / 0.805\\
\bottomrule\end{tabular}\end{adjustbox}\end{table}
\begin{table}[H]\centering\small
\caption*{Table S11 (continued). ETTh2}
\textbf{ETTh2}\par\smallskip
\begin{adjustbox}{max width=\linewidth}\begin{tabular}{lrrrrr}\toprule
\textbf{Model} & \textbf{96} & \textbf{192} & \textbf{336} & \textbf{720} & \textbf{Avg}\\\midrule
\rowcolor{blue!5}
\textbf{RDTU} & \textbf{0.263 / 0.329} & \textbf{0.334 / 0.372} & \textbf{0.355 / 0.381} & \textbf{0.395 / 0.425} & \textbf{0.337 / 0.377}\\
Autoformer & 0.346 / 0.388 & 0.456 / 0.452 & 0.482 / 0.486 & 0.515 / 0.511 & 0.450 / 0.459\\
Stationary & 0.476 / 0.458 & 0.512 / 0.493 & 0.552 / 0.551 & 0.562 / 0.560 & 0.526 / 0.516\\
ETSformer & 0.340 / 0.391 & 0.430 / 0.439 & 0.485 / 0.479 & 0.500 / 0.497 & 0.439 / 0.452\\
LightTS & 0.397 / 0.437 & 0.520 / 0.504 & 0.626 / 0.559 & 0.863 / 0.672 & 0.602 / 0.543\\
Informer & 3.755 / 1.525 & 5.602 / 1.931 & 4.721 / 1.835 & 3.647 / 1.625 & 4.431 / 1.729\\
Reformer & 2.626 / 1.317 & 11.12 / 2.979 & 9.323 / 2.769 & 3.874 / 1.697 & 6.736 / 2.191\\
\bottomrule\end{tabular}\end{adjustbox}\end{table}
\begin{table}[H]\centering\small
\caption*{Table S11 (continued). ETTm1}
\textbf{ETTm1}\par\smallskip
\begin{adjustbox}{max width=\linewidth}\begin{tabular}{lrrrrr}\toprule
\textbf{Model} & \textbf{96} & \textbf{192} & \textbf{336} & \textbf{720} & \textbf{Avg}\\\midrule
\rowcolor{blue!5}
\textbf{RDTU} & \textbf{0.272 / 0.326} & \textbf{0.327 / 0.354} & \textbf{0.350 / 0.375} & \textbf{0.410 / 0.412} & \textbf{0.340 / 0.367}\\
Autoformer & 0.505 / 0.475 & 0.553 / 0.496 & 0.621 / 0.537 & 0.671 / 0.561 & 0.588 / 0.517\\
Stationary & 0.386 / 0.398 & 0.459 / 0.444 & 0.495 / 0.464 & 0.585 / 0.516 & 0.481 / 0.456\\
ETSformer & 0.375 / 0.398 & 0.408 / 0.410 & 0.435 / 0.428 & 0.499 / 0.462 & 0.429 / 0.425\\
LightTS & 0.374 / 0.400 & 0.400 / 0.407 & 0.438 / 0.438 & 0.527 / 0.502 & 0.435 / 0.437\\
Informer & 0.672 / 0.571 & 0.795 / 0.669 & 1.212 / 0.871 & 1.166 / 0.823 & 0.961 / 0.734\\
Reformer & 0.538 / 0.528 & 0.658 / 0.592 & 0.898 / 0.721 & 1.102 / 0.841 & 0.799 / 0.671\\
\bottomrule\end{tabular}\end{adjustbox}\end{table}
\begin{table}[H]\centering\small
\caption*{Table S11 (continued). ETTm2}
\textbf{ETTm2}\par\smallskip
\begin{adjustbox}{max width=\linewidth}\begin{tabular}{lrrrrr}\toprule
\textbf{Model} & \textbf{96} & \textbf{192} & \textbf{336} & \textbf{720} & \textbf{Avg}\\\midrule
\rowcolor{blue!5}
\textbf{RDTU} & \textbf{0.157 / 0.244} & \textbf{0.219 / 0.288} & \textbf{0.270 / 0.323} & \textbf{0.354 / 0.379} & \textbf{0.250 / 0.309}\\
Autoformer & 0.255 / 0.339 & 0.281 / 0.340 & 0.339 / 0.372 & 0.433 / 0.432 & 0.327 / 0.371\\
Stationary & 0.192 / 0.274 & 0.280 / 0.339 & 0.334 / 0.361 & 0.417 / 0.413 & 0.306 / 0.347\\
ETSformer & 0.189 / 0.280 & 0.253 / 0.319 & 0.314 / 0.357 & 0.414 / 0.413 & 0.293 / 0.342\\
LightTS & 0.209 / 0.308 & 0.311 / 0.382 & 0.442 / 0.466 & 0.675 / 0.587 & 0.409 / 0.436\\
Informer & 0.365 / 0.453 & 0.533 / 0.563 & 1.363 / 0.887 & 3.379 / 1.338 & 1.410 / 0.810\\
Reformer & 0.658 / 0.619 & 1.078 / 0.827 & 1.549 / 0.972 & 2.631 / 1.242 & 1.479 / 0.915\\
\bottomrule\end{tabular}\end{adjustbox}\end{table}
\begin{table}[H]\centering\small
\caption*{Table S11 (continued). Weather}
\textbf{Weather}\par\smallskip
\begin{adjustbox}{max width=\linewidth}\begin{tabular}{lrrrrr}\toprule
\textbf{Model} & \textbf{96} & \textbf{192} & \textbf{336} & \textbf{720} & \textbf{Avg}\\\midrule
\rowcolor{blue!5}
\textbf{RDTU} & \textbf{0.141 / 0.184} & \textbf{0.189 / 0.227} & \textbf{0.239 / 0.272} & \textbf{0.314 / 0.322} & \textbf{0.221 / 0.251}\\
Autoformer & 0.266 / 0.336 & 0.307 / 0.367 & 0.359 / 0.395 & 0.419 / 0.428 & 0.338 / 0.382\\
Stationary & 0.173 / 0.223 & 0.245 / 0.285 & 0.321 / 0.338 & 0.414 / 0.410 & 0.288 / 0.314\\
ETSformer & 0.197 / 0.281 & 0.237 / 0.312 & 0.298 / 0.353 & 0.352 / 0.288 & 0.271 / 0.334\\
LightTS & 0.182 / 0.242 & 0.227 / 0.287 & 0.282 / 0.334 & 0.352 / 0.386 & 0.261 / 0.312\\
Informer & 0.300 / 0.384 & 0.598 / 0.544 & 0.578 / 0.523 & 1.059 / 0.741 & 0.634 / 0.548\\
Reformer & 0.689 / 0.596 & 0.752 / 0.638 & 0.639 / 0.596 & 1.130 / 0.792 & 0.803 / 0.656\\
\bottomrule\end{tabular}\end{adjustbox}\end{table}
\begin{table}[H]\centering\small
\caption*{Table S11 (continued). Electricity}
\textbf{Electricity}\par\smallskip
\begin{adjustbox}{max width=\linewidth}\begin{tabular}{lrrrrr}\toprule
\textbf{Model} & \textbf{96} & \textbf{192} & \textbf{336} & \textbf{720} & \textbf{Avg}\\\midrule
\rowcolor{blue!5}
\textbf{RDTU} & \textbf{0.124 / 0.210} & \textbf{0.144 / 0.235} & \textbf{0.162 / 0.257} & \textbf{0.193 / 0.281} & \textbf{0.156 / 0.246}\\
Autoformer & 0.201 / 0.317 & 0.222 / 0.334 & 0.231 / 0.338 & 0.254 / 0.361 & 0.227 / 0.338\\
Stationary & 0.169 / 0.273 & 0.182 / 0.286 & 0.200 / 0.304 & 0.222 / 0.321 & 0.193 / 0.296\\
ETSformer & 0.187 / 0.304 & 0.199 / 0.315 & 0.212 / 0.329 & 0.233 / 0.345 & 0.208 / 0.323\\
LightTS & 0.207 / 0.307 & 0.213 / 0.316 & 0.230 / 0.333 & 0.265 / 0.360 & 0.229 / 0.329\\
Informer & 0.274 / 0.368 & 0.296 / 0.386 & 0.300 / 0.394 & 0.373 / 0.439 & 0.311 / 0.397\\
Reformer & 0.312 / 0.402 & 0.348 / 0.433 & 0.350 / 0.433 & 0.340 / 0.420 & 0.338 / 0.422\\
\bottomrule\end{tabular}\end{adjustbox}\end{table}
\begin{table}[H]\centering\small
\caption*{Table S11 (continued). Traffic}
\textbf{Traffic}\par\smallskip
\begin{adjustbox}{max width=\linewidth}\begin{tabular}{lrrrrr}\toprule
\textbf{Model} & \textbf{96} & \textbf{192} & \textbf{336} & \textbf{720} & \textbf{Avg}\\\midrule
\rowcolor{blue!5}
\textbf{RDTU} & \textbf{0.358 / 0.235} & \textbf{0.375 / 0.248} & \textbf{0.388 / 0.256} & \textbf{0.435 / 0.280} & \textbf{0.389 / 0.255}\\
Autoformer & 0.613 / 0.388 & 0.616 / 0.382 & 0.622 / 0.337 & 0.660 / 0.408 & 0.628 / 0.379\\
Stationary & 0.612 / 0.338 & 0.613 / 0.340 & 0.618 / 0.328 & 0.653 / 0.355 & 0.624 / 0.340\\
ETSformer & 0.607 / 0.392 & 0.621 / 0.399 & 0.622 / 0.396 & 0.632 / 0.396 & 0.621 / 0.396\\
LightTS & 0.615 / 0.391 & 0.601 / 0.382 & 0.613 / 0.386 & 0.658 / 0.407 & 0.622 / 0.392\\
Informer & 0.719 / 0.391 & 0.696 / 0.379 & 0.777 / 0.420 & 0.864 / 0.472 & 0.764 / 0.416\\
Reformer & 0.732 / 0.423 & 0.733 / 0.420 & 0.742 / 0.420 & 0.755 / 0.423 & 0.741 / 0.422\\
\bottomrule\end{tabular}\end{adjustbox}\end{table}
\begin{table}[H]\centering\small
\caption*{Table S11 (continued). ILI}
\textbf{ILI}\par\smallskip
\begin{adjustbox}{max width=\linewidth}\begin{tabular}{lrrrrr}\toprule
\textbf{Model} & \textbf{24} & \textbf{36} & \textbf{48} & \textbf{60} & \textbf{Avg}\\\midrule
\rowcolor{blue!5}
\textbf{RDTU} & \textbf{1.609 / 0.771} & \textbf{1.453 / 0.835} & \textbf{1.545 / 0.841} & \textbf{1.630 / 0.802} & \textbf{1.559 / 0.812}\\
Autoformer & 3.483 / 1.287 & 3.103 / 1.148 & 2.669 / 1.085 & 2.770 / 1.125 & 3.006 / 1.161\\
Stationary & 2.294 / 0.945 & 1.825 / 0.848 & 2.010 / 0.900 & 2.178 / 0.963 & 2.077 / 0.914\\
ETSformer & 2.527 / 1.020 & 2.615 / 1.007 & 2.359 / 0.972 & 2.487 / 1.016 & 2.497 / 1.004\\
LightTS & 8.313 / 2.144 & 6.631 / 1.902 & 7.299 / 1.982 & 7.283 / 1.985 & 7.382 / 2.003\\
Informer & 5.764 / 1.677 & 4.755 / 1.467 & 4.763 / 1.469 & 5.264 / 1.564 & 5.137 / 1.544\\
Reformer & 4.400 / 1.382 & 4.783 / 1.448 & 4.832 / 1.465 & 4.882 / 1.483 & 4.724 / 1.445\\
\bottomrule\end{tabular}\end{adjustbox}\end{table}
\noindent\textbf{Source-reported ranking counts.} RDTU: 40, Autoformer: 0, Stationary: 0, ETSformer: 0, LightTS: 0, Informer: 0, Reformer: 0. The original counting convention is not specified; these values are retained without recalculation.\par



\shortautoref{tab:long-term-forecasting-additional} presents an extensive evaluation of the proposed RDTU method against six state-of-the-art baseline models (Autoformer, Stationary, ETSformer, LightTS, Informer, and Reformer) across eight benchmark datasets. The experimental results clearly demonstrate the superior performance of RDTU for long-term time series forecasting.

As indicated by the bold red values, RDTU achieves the lowest Mean Squared Error (MSE) across the listed settings and the lowest Mean Absolute Error (MAE) in all but one individual horizon setting ($H \in \{24, 36, 48, 60\}$ for ILI and $H \in \{96, 192, 336, 720\}$ for other datasets). The quantitative summary at the bottom of the table confirms this dominance, with a source-reported "1st Count" of 40 for RDTU and 0 for each listed baseline. This count is retained as reported; it is distinct from the number of individual MSE and MAE cells in the table. While models such as LightTS and ETSformer occasionally yield the second-best results (highlighted in blue), RDTU consistently provides significant accuracy improvements.


\subsection{Short-term Forecasting}

Our complete results on short-term forecasting are presented in \shortautoref{tab:short-term-forecasting}.

\shortautoref{tab:short-term-forecasting} extends the analysis to short-term time series forecasting. Evaluation metrics including SMAPE, MASE, and OWA indicate that RDTU maintains its leading position. In the weighted average analysis across all sampling intervals (Yearly, Quarterly, Monthly, Others), RDTU records the best performance (e.g., Average OWA of 0.859), and matches Time-LLM on OWA while achieving lower SMAPE (OWA 0.859), significantly surpassing traditional methods like N-HiTS and N-BEATS.

\begin{table}[H]\centering\small
\caption{Full M4 short-term forecasting results. Lower values are better.}\label{tab:short-term-forecasting}
\textbf{Yearly}\par\smallskip
\begin{adjustbox}{max width=\linewidth}\begin{tabular}{lrrr}\toprule
\textbf{Model} & \textbf{SMAPE} & \textbf{MASE} & \textbf{OWA}\\\midrule
\rowcolor{blue!5}
\textbf{RDTU} & \textbf{13.408} & \textbf{3.014} & \textbf{0.788}\\
Time-LLM & 13.419 & 3.005 & 0.789\\
GPT4TS & 15.11 & 3.565 & 0.911\\
TimesNet & 15.378 & 3.554 & 0.918\\
PatchTST & 13.477 & 3.019 & 0.792\\
N-HiTS & 13.422 & 3.056 & 0.795\\
N-BEATS & 13.487 & 3.036 & 0.795\\
ETSformer & 18.009 & 4.487 & 1.115\\
LightTS & 14.247 & 3.109 & 0.827\\
DLinear & 16.965 & 4.283 & 1.058\\
FEDformer & 14.021 & 3.036 & 0.811\\
Stationary & 13.717 & 3.078 & 0.807\\
Autoformer & 13.974 & 3.134 & 0.822\\
Informer & 14.727 & 3.418 & 0.881\\
Reformer & 16.169 & 3.800 & 0.973\\
\bottomrule\end{tabular}\end{adjustbox}\end{table}
\begin{table}[H]\centering\small
\caption*{Table S12 (continued). Quarterly}
\textbf{Quarterly}\par\smallskip
\begin{adjustbox}{max width=\linewidth}\begin{tabular}{lrrr}\toprule
\textbf{Model} & \textbf{SMAPE} & \textbf{MASE} & \textbf{OWA}\\\midrule
\rowcolor{blue!5}
\textbf{RDTU} & \textbf{10.094} & \textbf{1.176} & \textbf{0.888}\\
Time-LLM & 10.110 & 1.178 & 0.889\\
GPT4TS & 10.597 & 1.253 & 0.938\\
TimesNet & 10.465 & 1.227 & 0.923\\
PatchTST & 10.38 & 1.233 & 0.921\\
N-HiTS & 10.185 & 1.18 & 0.893\\
N-BEATS & 10.564 & 1.252 & 0.936\\
ETSformer & 13.376 & 1.906 & 1.302\\
LightTS & 11.364 & 1.328 & 1.000\\
DLinear & 12.145 & 1.520 & 1.106\\
FEDformer & 11.1 & 1.35 & 0.996\\
Stationary & 10.958 & 1.325 & 0.981\\
Autoformer & 11.338 & 1.365 & 1.012\\
Informer & 11.360 & 1.401 & 1.027\\
Reformer & 13.313 & 1.775 & 1.252\\
\bottomrule\end{tabular}\end{adjustbox}\end{table}
\begin{table}[H]\centering\small
\caption*{Table S12 (continued). Monthly}
\textbf{Monthly}\par\smallskip
\begin{adjustbox}{max width=\linewidth}\begin{tabular}{lrrr}\toprule
\textbf{Model} & \textbf{SMAPE} & \textbf{MASE} & \textbf{OWA}\\\midrule
\rowcolor{blue!5}
\textbf{RDTU} & \textbf{12.977} & \textbf{0.968} & \textbf{0.903}\\
Time-LLM & 12.980 & 0.963 & 0.903\\
GPT4TS & 13.258 & 1.003 & 0.931\\
TimesNet & 13.513 & 1.039 & 0.957\\
PatchTST & 12.959 & 0.97 & 0.905\\
N-HiTS & 13.059 & 1.013 & 0.929\\
N-BEATS & 13.089 & 0.996 & 0.922\\
ETSformer & 14.588 & 1.368 & 1.149\\
LightTS & 14.014 & 1.053 & 0.981\\
DLinear & 13.514 & 1.037 & 0.956\\
FEDformer & 14.403 & 1.147 & 1.038\\
Stationary & 13.917 & 1.097 & 0.998\\
Autoformer & 13.958 & 1.103 & 1.002\\
Informer & 14.062 & 1.141 & 1.024\\
Reformer & 20.128 & 2.614 & 1.927\\
\bottomrule\end{tabular}\end{adjustbox}\end{table}
\begin{table}[H]\centering\small
\caption*{Table S12 (continued). Others}
\textbf{Others}\par\smallskip
\begin{adjustbox}{max width=\linewidth}\begin{tabular}{lrrr}\toprule
\textbf{Model} & \textbf{SMAPE} & \textbf{MASE} & \textbf{OWA}\\\midrule
\rowcolor{blue!5}
\textbf{RDTU} & \textbf{4.730} & \textbf{3.086} & \textbf{0.985}\\
Time-LLM & 4.795 & 3.178 & 1.006\\
GPT4TS & 6.124 & 4.116 & 1.259\\
TimesNet & 6.913 & 4.507 & 1.438\\
PatchTST & 4.952 & 3.347 & 1.049\\
N-HiTS & 4.711 & 3.054 & 0.977\\
N-BEATS & 6.599 & 4.43 & 1.393\\
ETSformer & 7.267 & 5.240 & 1.591\\
LightTS & 15.880 & 11.434 & 3.474\\
DLinear & 6.709 & 4.953 & 1.487\\
FEDformer & 7.148 & 4.041 & 1.389\\
Stationary & 6.302 & 4.064 & 1.304\\
Autoformer & 5.485 & 3.865 & 1.187\\
Informer & 24.460 & 20.960 & 5.879\\
Reformer & 32.491 & 33.355 & 8.679\\
\bottomrule\end{tabular}\end{adjustbox}\end{table}
\begin{table}[H]\centering\small
\caption*{Table S12 (continued). Average}
\textbf{Average}\par\smallskip
\begin{adjustbox}{max width=\linewidth}\begin{tabular}{lrrr}\toprule
\textbf{Model} & \textbf{SMAPE} & \textbf{MASE} & \textbf{OWA}\\\midrule
\rowcolor{blue!5}
\textbf{RDTU} & \textbf{11.979} & \textbf{1.599} & \textbf{0.859}\\
Time-LLM & 11.983 & 1.595 & 0.859\\
GPT4TS & 12.69 & 1.808 & 0.94\\
TimesNet & 12.88 & 1.836 & 0.955\\
PatchTST & 12.059 & 1.623 & 0.869\\
N-HiTS & 12.035 & 1.625 & 0.869\\
N-BEATS & 12.25 & 1.698 & 0.896\\
ETSformer & 14.718 & 2.408 & 1.172\\
LightTS & 13.525 & 2.111 & 1.051\\
DLinear & 13.639 & 2.095 & 1.051\\
FEDformer & 13.16 & 1.775 & 0.949\\
Stationary & 12.780 & 1.756 & 0.930\\
Autoformer & 12.909 & 1.771 & 0.939\\
Informer & 14.086 & 2.718 & 1.230\\
Reformer & 18.200 & 4.223 & 1.775\\
\bottomrule\end{tabular}\end{adjustbox}\end{table}

